\chapter{Tema, Problema, Objetivos e Justificativa}

A escolha do tema é uma das decisões mais importantes do TCC. Um tema adequado precisa ser relevante, viável, delimitado e compatível com a área de formação. A partir dele, o estudante constrói problema, objetivos e justificativa.

\section{Escolha do tema}

O tema pode surgir de experiências profissionais, disciplinas cursadas, estágios, projetos de extensão, iniciação científica, observação de problemas sociais, lacunas na literatura ou demandas institucionais. O interesse pessoal é importante, mas não basta. O tema precisa permitir investigação acadêmica.

Antes de decidir, verifique:

\begin{itemize}
  \item se há bibliografia suficiente e atual;
  \item se o recorte cabe no tempo disponível;
  \item se há acesso aos dados, documentos, campo ou participantes;
  \item se o orientador possui aderência ao tema;
  \item se há necessidade de aprovação ética;
  \item se o tema dialoga com a formação do curso.
\end{itemize}

\section{Delimitação}

Delimitar é transformar um assunto amplo em objeto investigável. Um tema como ``saúde mental de estudantes'' é amplo demais. Um recorte mais viável poderia considerar um curso, uma instituição, um período, um grupo específico, uma intervenção, uma política ou uma dimensão do fenômeno.

A delimitação pode envolver:

\begin{itemize}
  \item espaço: instituição, território, serviço, comunidade ou ambiente;
  \item tempo: período histórico, ano, semestre, ciclo ou marco específico;
  \item população: grupo, perfil, categoria profissional ou público-alvo;
  \item material: documentos, bases, obras, sistemas, processos ou práticas;
  \item enfoque: dimensão teórica, metodológica, política, técnica ou social.
\end{itemize}

\section{Problema de pesquisa}

O problema de pesquisa é a pergunta central do trabalho. Ele deve ser claro, investigável e coerente com o método. Um problema muito amplo leva a objetivos genéricos. Um problema muito restrito pode impedir análise significativa.

Perguntas de pesquisa costumam usar verbos e expressões como: de que modo, em que medida, quais fatores, quais efeitos, como ocorre, quais relações, que percepções, que desafios, que contribuições.

\caixaorientacao{Um bom teste é perguntar: ao final do trabalho, será possível responder a essa pergunta com base no método escolhido? Se a resposta for não, o problema ainda precisa ser ajustado.}

\section{Objetivo geral e objetivos específicos}

O objetivo geral expressa a finalidade principal do trabalho. Os objetivos específicos indicam etapas necessárias para alcançar o objetivo geral. Eles não devem ser uma lista de tarefas mecânicas, mas ações intelectuais ligadas à investigação.

Verbos úteis incluem analisar, compreender, descrever, identificar, comparar, avaliar, investigar, discutir, mapear, propor, sistematizar e interpretar. Verbos vagos, como conhecer ou falar sobre, devem ser evitados.

\section{Justificativa}

A justificativa explica por que o trabalho deve ser realizado. Ela pode apresentar relevância acadêmica, social, institucional, profissional, científica ou tecnológica. Também pode indicar lacunas, demandas, problemas práticos ou necessidade de sistematização.

Uma justificativa consistente não depende de afirmações grandiosas. Ela demonstra, com argumentos e fontes, que o tema merece atenção e que o recorte proposto é pertinente.

